\section{Discussion and Limitations}
\label{sec:discussion}

The selected system will run as a standalone application on
the university's Linux machines, without internet resources
or existing AI systems.

\paragraph{Search effort.}
Finding an initial valid timetable may require exploring many
assignments, particularly when capacities and other hard
constraints leave only a few choices. If no valid timetable exists,
the OR-tree search must exhaust its alternatives before
reporting failure. Choosing the component with the fewest
legal slots helps expose alternatives that do not work early, but does not
guarantee a short search.

\paragraph{Limits of the optimization stages.}
Set-based search accepts only moves and swaps that immediately
reduce the penalty. It can therefore stop even when a better
timetable exists, because reaching that timetable may require
changing several components together or first accepting a
change that does not improve the score.

AND-tree search can establish an optimal result, but may need
to examine many cases. Our lower bound counts only preference
penalties already fixed by the current assignments. This makes
the bound straightforward to check, but it ignores other
penalties that may be unavoidable. As a result, some cases
will remain in the search even though they cannot improve
the best timetable. If all preference penalties are zero,
this bound provides no pruning while the best known penalty
is positive.

\paragraph{Time limits.}
We plan to use time limits in our implementation to control running
time on larger instances. If the OR-tree search reaches this
limit before finding a valid timetable, it will report that
the search stopped without a solution, which does not establish that no valid timetable exists. If the limit is reached during
either optimization stage, the system will return the best
valid timetable found so far. It will report whether the search
finished or stopped at the time limit, and will not claim
optimality unless it has been established. A timetable with
penalty zero is optimal because all penalties are nonnegative.

\paragraph{Verification.}
We will test each hard constraint and penalty calculation
on small inputs with known results. On small 
instances, an exhaustive search will provide the minimum
penalty for comparison. Proposal B should match that minimum
when its search finishes; Proposal A should return a valid
timetable no worse than its initial timetable, though it may
not reach the minimum. Larger inputs will be used to assess
running time and solution quality.